%% Generated by Sphinx.
\def\sphinxdocclass{report}
\documentclass[letterpaper,10pt,english]{sphinxmanual}
\ifdefined\pdfpxdimen
   \let\sphinxpxdimen\pdfpxdimen\else\newdimen\sphinxpxdimen
\fi \sphinxpxdimen=.75bp\relax

\PassOptionsToPackage{warn}{textcomp}
\usepackage[utf8]{inputenc}
\ifdefined\DeclareUnicodeCharacter
% support both utf8 and utf8x syntaxes
  \ifdefined\DeclareUnicodeCharacterAsOptional
    \def\sphinxDUC#1{\DeclareUnicodeCharacter{"#1}}
  \else
    \let\sphinxDUC\DeclareUnicodeCharacter
  \fi
  \sphinxDUC{00A0}{\nobreakspace}
  \sphinxDUC{2500}{\sphinxunichar{2500}}
  \sphinxDUC{2502}{\sphinxunichar{2502}}
  \sphinxDUC{2514}{\sphinxunichar{2514}}
  \sphinxDUC{251C}{\sphinxunichar{251C}}
  \sphinxDUC{2572}{\textbackslash}
\fi
\usepackage{cmap}
\usepackage[T1]{fontenc}
\usepackage{amsmath,amssymb,amstext}
\usepackage{babel}



\usepackage{times}
\expandafter\ifx\csname T@LGR\endcsname\relax
\else
% LGR was declared as font encoding
  \substitutefont{LGR}{\rmdefault}{cmr}
  \substitutefont{LGR}{\sfdefault}{cmss}
  \substitutefont{LGR}{\ttdefault}{cmtt}
\fi
\expandafter\ifx\csname T@X2\endcsname\relax
  \expandafter\ifx\csname T@T2A\endcsname\relax
  \else
  % T2A was declared as font encoding
    \substitutefont{T2A}{\rmdefault}{cmr}
    \substitutefont{T2A}{\sfdefault}{cmss}
    \substitutefont{T2A}{\ttdefault}{cmtt}
  \fi
\else
% X2 was declared as font encoding
  \substitutefont{X2}{\rmdefault}{cmr}
  \substitutefont{X2}{\sfdefault}{cmss}
  \substitutefont{X2}{\ttdefault}{cmtt}
\fi


\usepackage[Bjarne]{fncychap}
\usepackage[,numfigreset=1,mathnumfig]{sphinx}

\fvset{fontsize=\small}
\usepackage{geometry}


% Include hyperref last.
\usepackage{hyperref}
% Fix anchor placement for figures with captions.
\usepackage{hypcap}% it must be loaded after hyperref.
% Set up styles of URL: it should be placed after hyperref.
\urlstyle{same}

\addto\captionsenglish{\renewcommand{\contentsname}{Contents:}}

\usepackage{sphinxmessages}
\setcounter{tocdepth}{1}



\title{Agrogeophy-catalog Documentation}
\date{Feb 22, 2021}
\release{0.1.0}
\author{agrogeophy}
\newcommand{\sphinxlogo}{\vbox{}}
\renewcommand{\releasename}{Release}
\makeindex
\begin{document}

\pagestyle{empty}
\sphinxmaketitle
\pagestyle{plain}
\sphinxtableofcontents
\pagestyle{normal}
\phantomsection\label{\detokenize{index::doc}}



\chapter{Getting Started}
\label{\detokenize{getting-started:getting-started}}\label{\detokenize{getting-started::doc}}
\sphinxhref{https://agrogeophy.slack.com/}{\sphinxincludegraphics{{/home/runner/work/catalog/catalog/docs/tex/.doctrees/images/9ca7c69c7a888bdf4340e21de765bed9e1a8b192/Slack-agrogeophy-1}.svg}}

\sphinxhref{https://doi.org/10.5281/zenodo.4058524}{\sphinxincludegraphics{{/home/runner/work/catalog/catalog/docs/tex/.doctrees/images/0a96cfbd8573408ea7e006027cb81d65ecfa3622/zenodo.4058524}.svg}}


\section{About}
\label{\detokenize{getting-started:about}}
\sphinxAtStartPar
\sphinxstylestrong{Agrogeophy\sphinxhyphen{}catalog is part of the agrogeophysical catalog website:} \sphinxurl{http://geo.geoscienze.unipd.it/growingwebsite/map\_catalog}

\sphinxAtStartPar
We welcome any feedback and ideas!
Let us know by submitting
\sphinxhref{https://github.com/BenjMy/agrogeophy-catalog/issues}{issues on Github}
or send us a message on our
\sphinxhref{https://agrogeophy.slack.com/}{Slack chatroom}.


\section{Citing}
\label{\detokenize{getting-started:citing}}
\sphinxAtStartPar
If you use agrogeophy\sphinxhyphen{}catalog for you work, please cite this paper as:

\sphinxAtStartPar
TODO

\sphinxAtStartPar
BibTex code:

\begin{sphinxVerbatim}[commandchars=\\\{\}]
\PYG{n}{TODO}
\end{sphinxVerbatim}


\chapter{CAGS ecosystem}
\label{\detokenize{CAGS_ecosystem:cags-ecosystem}}\label{\detokenize{CAGS_ecosystem::doc}}
\sphinxAtStartPar
CAGS is an ecosystem with many interrelated projects and repositories:
\begin{itemize}
\item {} 
\sphinxAtStartPar
\sphinxstylestrong{Catalog}: a platform putting together a database/catalog of agrogeophysical surveys in order to boost future research

\item {} 
\sphinxAtStartPar
\sphinxstylestrong{Notebooks}: a collection of jupyter notebooks that use agrogeophysical models

\end{itemize}

\sphinxAtStartPar
\sphinxhref{https://agrogeophy.github.io/notebooks/}{Documentation}
\begin{itemize}
\item {} 
\sphinxAtStartPar
\sphinxstylestrong{Datasets}: a list of agrogeophysical dataset of interest with associated metadata and references

\end{itemize}

\sphinxAtStartPar
\sphinxhref{https://agrogeophy.github.io/datasets/}{Documentation}

\begin{figure}[htbp]
\centering
\capstart

\noindent\sphinxincludegraphics[scale=0.6]{{CAGS_ecosystem}.png}
\caption{Scheme describing the ecosystem of the CAGS collaborative platform and the different paths for user’s contribution. (Path A) Article and/or companion notebook submission to the catalogue (B.1) optional companion dataset (B.2) companion dataset complying with CAGS preparation guidelines. (C) the data preparation for a submission to the CAGS datasets for key parameters described in section 4.3). (D) all possible interactions and management of submissions. The CAGS dataset preparation guidelines are described in section 2.2. Two metadata harvesters are presented i.e. the project\sphinxhyphen{}level and the file\sphinxhyphen{}level geophysics one}\label{\detokenize{CAGS_ecosystem:id2}}\label{\detokenize{CAGS_ecosystem:importing}}\end{figure}


\chapter{CAGS \& FAIR principles}
\label{\detokenize{FAIR:cags-fair-principles}}\label{\detokenize{FAIR::doc}}
\sphinxAtStartPar
FAIR: Findability, Accessibility, Interoperability, and Reuse


\section{Findability}
\label{\detokenize{FAIR:findability}}
\sphinxAtStartPar
To enable an effective findability of users contributions, the website interface automatically extracts and inserts into the database all the metadata provided in the user submission form (for each type i.e. datasets, numerical studies and communications). The website interface allows then the retrieval and display of the contribution thanks to its search and filtering tool. A geospatial visualization and analysis component enables searching, visualizing, and analyzing users’ contributions.

\begin{figure}[htbp]
\centering
\capstart

\noindent\sphinxincludegraphics[scale=0.5]{{Map_contributions_26102020}.png}
\caption{Map of georeferenced studies (zoom on the Mediterranean basin). The map is based on the Leaflet framework and use map background from the OpenStreetMap project.}\label{\detokenize{FAIR:id1}}\label{\detokenize{FAIR:importing}}\end{figure}


\section{Accessibility}
\label{\detokenize{FAIR:accessibility}}
\sphinxAtStartPar
Once the user finds the required data, it can be access thanks to its DOI link. Metadata are accessible, even when the data are no longer available.


\section{Interoperability}
\label{\detokenize{FAIR:interoperability}}
\sphinxAtStartPar
The concept of standardized descriptive metadata provides a powerful mechanism to improve retrieval for specific applications and user communities and facilitates repository interoperability. The catalogue metadata architecture we implemented for CGAS has been inspired by the Archaeology Data Service (ADS, Richards, 2018) which acts as a metadata aggregator between archaeological and geophysical metadata.


\section{Reuse}
\label{\detokenize{FAIR:reuse}}
\sphinxAtStartPar
As a start, CAGS implement useful tools for the dataset:
\begin{quote}

\sphinxAtStartPar
To make this decision, the data publisher should provide not just metadata that allows discovery, but also metadata that richly describes the context under which the data was generated.
\end{quote}
\begin{itemize}
\item {} 
\sphinxAtStartPar
A data package linter , which could check for structure and files

\item {} 
\sphinxAtStartPar
At the raw data level, REDA may also serve as a translation service, converting old formats to new ones when necessary, especially for IP/SIP data, a fact that will highly contribute to the long\sphinxhyphen{}term data management and reuse.

\item {} 
\sphinxAtStartPar
At the (pre)\sphinxhyphen{}processed data level, the user may want to share the journal file produced using the REDA Python package.

\end{itemize}


\chapter{Glossary}
\label{\detokenize{glossary:glossary}}\label{\detokenize{glossary::doc}}

\section{Geoelectrical methods}
\label{\detokenize{glossary:geoelectrical-methods}}

\subsection{Graphical}
\label{\detokenize{glossary:graphical}}
\begin{figure}[htbp]
\centering
\capstart

\noindent\sphinxincludegraphics[scale=0.1]{{channels}.png}
\caption{Possible channels for soil/plant geoelectical measurements. Credit to L. Peruzzo {[}CIT2002{]}.}\label{\detokenize{glossary:id7}}\label{\detokenize{glossary:importing}}\end{figure}

\begin{figure}[htbp]
\centering

\noindent\sphinxincludegraphics[scale=0.5]{{ehosioke_sensing_2020_fig2}.png}
\end{figure}

\sphinxAtStartPar
See \sphinxcite{glossary:id6} for more details.

\sphinxAtStartPar
Figure 2


\subsection{Acronyms}
\label{\detokenize{glossary:acronyms}}
\sphinxAtStartPar
\sphinxstylestrong{spectroscopic}
\begin{itemize}
\item {} 
\sphinxAtStartPar
EIS: Electrical Impedance Spectroscopy

\item {} 
\sphinxAtStartPar
IP: Induced Polarization

\item {} 
\sphinxAtStartPar
SIP: Spectral Induced Polarization

\item {} 
\sphinxAtStartPar
TDIP: Time\sphinxhyphen{}domain Induced Polarization

\end{itemize}

\sphinxAtStartPar
\sphinxstylestrong{imaging or tomography}
\begin{itemize}
\item {} 
\sphinxAtStartPar
ERT: Electrical Resistivity Tomography

\item {} 
\sphinxAtStartPar
ERI: Electrical Resistivity Imaging

\item {} 
\sphinxAtStartPar
ECT: Electrical Capacitance Tomography

\item {} 
\sphinxAtStartPar
EIT: Electrical Impedance Tomography

\item {} 
\sphinxAtStartPar
sEIT: spectral Electrical Impedance Tomography

\end{itemize}


\section{EM methods}
\label{\detokenize{glossary:em-methods}}

\subsection{Acronyms}
\label{\detokenize{glossary:id4}}\begin{itemize}
\item {} 
\sphinxAtStartPar
EMI: ElectroMagnetic Induction

\end{itemize}


\section{References}
\label{\detokenize{glossary:references}}
\sphinxAtStartPar



\chapter{The CAGS Metadata Schema}
\label{\detokenize{schema_documentation:the-cags-metadata-schema}}\label{\detokenize{schema_documentation::doc}}

\section{Project level metadata}
\label{\detokenize{schema_documentation:project-level-metadata}}
\sphinxAtStartPar
The \sphinxstylestrong{Project\sphinxhyphen{}level Metadata} level is related to the main CAGS interface website from which metadata are harvested through the online submission form. It collects metadata (derived and expended from Dublin Core metadata) which incorporates a number of descriptive and resource discovery focused elements describing either a scientific communication, a notebook and a dataset independently or jointly.

\sphinxAtStartPar
Tables below are extracted from the database and generated by MySQL Workbench Model Documentation v1.0.0 \sphinxhyphen{} Copyright (c)
2015 Hieu Le


\subsection{Table: \sphinxstyleliteralintitle{\sphinxupquote{abiotic}}}
\label{\detokenize{schema_documentation:table-abiotic}}

\subsubsection{Description:}
\label{\detokenize{schema_documentation:description}}

\subsubsection{Columns:}
\label{\detokenize{schema_documentation:columns}}

\begin{savenotes}\sphinxattablestart
\centering
\begin{tabulary}{\linewidth}[t]{|T|T|T|T|T|T|T|}
\hline
\sphinxstyletheadfamily 
\sphinxAtStartPar
Column
&\sphinxstyletheadfamily 
\sphinxAtStartPar
Description
&\sphinxstyletheadfamily 
\sphinxAtStartPar
Multiplicity
&\sphinxstyletheadfamily 
\sphinxAtStartPar
Obligation
&\sphinxstyletheadfamily 
\sphinxAtStartPar
Metadata level
&\sphinxstyletheadfamily 
\sphinxAtStartPar
Dublin Core
compatibility
&\sphinxstyletheadfamily 
\sphinxAtStartPar
INSPIRE
compatibility
\\
\hline
\sphinxAtStartPar
soil
&&
\sphinxAtStartPar
{[}1,n{]}
&
\sphinxAtStartPar
Optional
&
\sphinxAtStartPar
project
&
\sphinxAtStartPar
NONE
&
\sphinxAtStartPar
NONE
\\
\hline
\sphinxAtStartPar
water
&
\sphinxAtStartPar
Triggered infiltration test:    Yes/No
&
\sphinxAtStartPar
{[}1,n{]}
&
\sphinxAtStartPar
Optional
&
\sphinxAtStartPar
project
&
\sphinxAtStartPar
NONE
&
\sphinxAtStartPar
NONE
\\
\hline
\sphinxAtStartPar
land\_use
&&
\sphinxAtStartPar
{[}1,n{]}
&
\sphinxAtStartPar
Optional
&
\sphinxAtStartPar
project
&
\sphinxAtStartPar
NONE
&
\sphinxAtStartPar
NONE
\\
\hline
\end{tabulary}
\par
\sphinxattableend\end{savenotes}


\subsection{Table: \sphinxstyleliteralintitle{\sphinxupquote{biotic}}}
\label{\detokenize{schema_documentation:table-biotic}}

\subsubsection{Description:}
\label{\detokenize{schema_documentation:id1}}

\subsubsection{Columns:}
\label{\detokenize{schema_documentation:id2}}

\begin{savenotes}\sphinxattablestart
\centering
\begin{tabulary}{\linewidth}[t]{|T|T|T|T|T|T|T|}
\hline
\sphinxstyletheadfamily 
\sphinxAtStartPar
Column
&\sphinxstyletheadfamily 
\sphinxAtStartPar
Description
&\sphinxstyletheadfamily 
\sphinxAtStartPar
Multiplicity
&\sphinxstyletheadfamily 
\sphinxAtStartPar
Obligation
&\sphinxstyletheadfamily 
\sphinxAtStartPar
Metadata level
&\sphinxstyletheadfamily 
\sphinxAtStartPar
Dublin Core
compatibility
&\sphinxstyletheadfamily 
\sphinxAtStartPar
INSPIRE
compatibility
\\
\hline
\sphinxAtStartPar
species
&&
\sphinxAtStartPar
{[}1,n{]}
&
\sphinxAtStartPar
Optional
&
\sphinxAtStartPar
project
&
\sphinxAtStartPar
NONE
&
\sphinxAtStartPar
NONE
\\
\hline
\sphinxAtStartPar
organ
&&
\sphinxAtStartPar
{[}1,n{]}
&
\sphinxAtStartPar
Optional
&
\sphinxAtStartPar
project
&
\sphinxAtStartPar
NONE
&
\sphinxAtStartPar
NONE
\\
\hline
\end{tabulary}
\par
\sphinxattableend\end{savenotes}


\subsection{Table: \sphinxstyleliteralintitle{\sphinxupquote{contact}}}
\label{\detokenize{schema_documentation:table-contact}}

\subsubsection{Description:}
\label{\detokenize{schema_documentation:id3}}

\subsubsection{Columns:}
\label{\detokenize{schema_documentation:id4}}

\begin{savenotes}\sphinxattablestart
\centering
\begin{tabulary}{\linewidth}[t]{|T|T|T|T|T|T|T|}
\hline
\sphinxstyletheadfamily 
\sphinxAtStartPar
Column
&\sphinxstyletheadfamily 
\sphinxAtStartPar
Description
&\sphinxstyletheadfamily 
\sphinxAtStartPar
Multiplicity
&\sphinxstyletheadfamily 
\sphinxAtStartPar
Obligation
&\sphinxstyletheadfamily 
\sphinxAtStartPar
Metadata level
&\sphinxstyletheadfamily 
\sphinxAtStartPar
Dublin Core
compatibility
&\sphinxstyletheadfamily 
\sphinxAtStartPar
INSPIRE
compatibility
\\
\hline
\sphinxAtStartPar
name
&&
\sphinxAtStartPar
{[}1{]}
&
\sphinxAtStartPar
Mandatory
&
\sphinxAtStartPar
project
&
\sphinxAtStartPar
NONE
&
\sphinxAtStartPar
NONE
\\
\hline
\sphinxAtStartPar
surname
&&
\sphinxAtStartPar
{[}1{]}
&
\sphinxAtStartPar
Mandatory
&
\sphinxAtStartPar
project
&
\sphinxAtStartPar
NONE
&
\sphinxAtStartPar
NONE
\\
\hline
\sphinxAtStartPar
organisation
&&
\sphinxAtStartPar
{[}1,n{]}
&
\sphinxAtStartPar
Optional
&
\sphinxAtStartPar
project
&
\sphinxAtStartPar
NONE
&
\sphinxAtStartPar
NONE
\\
\hline
\sphinxAtStartPar
email
&&&
\sphinxAtStartPar
Mandatory
&
\sphinxAtStartPar
project
&
\sphinxAtStartPar
NONE
&
\sphinxAtStartPar
NONE
\\
\hline
\end{tabulary}
\par
\sphinxattableend\end{savenotes}


\subsection{Table: \sphinxstyleliteralintitle{\sphinxupquote{main}}}
\label{\detokenize{schema_documentation:table-main}}

\subsubsection{Description:}
\label{\detokenize{schema_documentation:id5}}

\subsubsection{Columns:}
\label{\detokenize{schema_documentation:id6}}

\begin{savenotes}\sphinxattablestart
\centering
\begin{tabulary}{\linewidth}[t]{|T|T|T|T|T|T|T|}
\hline
\sphinxstyletheadfamily 
\sphinxAtStartPar
Column
&\sphinxstyletheadfamily 
\sphinxAtStartPar
Description
&\sphinxstyletheadfamily 
\sphinxAtStartPar
Multiplicity
&\sphinxstyletheadfamily 
\sphinxAtStartPar
Obligation
&\sphinxstyletheadfamily 
\sphinxAtStartPar
Metadata level
&\sphinxstyletheadfamily 
\sphinxAtStartPar
Dublin Core
compatibility
&\sphinxstyletheadfamily 
\sphinxAtStartPar
INSPIRE
compatibility
\\
\hline
\sphinxAtStartPar
contrib\_type
&&
\sphinxAtStartPar
{[}1{]}
&
\sphinxAtStartPar
Mandatory
&
\sphinxAtStartPar
project
&
\sphinxAtStartPar
NONE
&
\sphinxAtStartPar
NONE
\\
\hline
\sphinxAtStartPar
contrib\_title
&&
\sphinxAtStartPar
{[}1{]}
&
\sphinxAtStartPar
Mandatory
&
\sphinxAtStartPar
project
&
\sphinxAtStartPar
NONE
&
\sphinxAtStartPar
NONE
\\
\hline
\sphinxAtStartPar
contrib\_date
&&
\sphinxAtStartPar
{[}1,n{]}
&
\sphinxAtStartPar
Optional
&
\sphinxAtStartPar
project
&
\sphinxAtStartPar
NONE
&
\sphinxAtStartPar
NONE
\\
\hline
\sphinxAtStartPar
contrib\_authors
&&&
\sphinxAtStartPar
Mandatory
&
\sphinxAtStartPar
project
&
\sphinxAtStartPar
NONE
&
\sphinxAtStartPar
NONE
\\
\hline
\sphinxAtStartPar
DOI
&&&
\sphinxAtStartPar
Mandatory
&
\sphinxAtStartPar
project
&
\sphinxAtStartPar
NONE
&
\sphinxAtStartPar
NONE
\\
\hline
\sphinxAtStartPar
journal
&&&
\sphinxAtStartPar
Optional
&
\sphinxAtStartPar
project
&
\sphinxAtStartPar
NONE
&
\sphinxAtStartPar
NONE
\\
\hline
\sphinxAtStartPar
icon\_img
&&&
\sphinxAtStartPar
Optional
&
\sphinxAtStartPar
project
&
\sphinxAtStartPar
NONE
&
\sphinxAtStartPar
NONE
\\
\hline
\sphinxAtStartPar
keywords
&&&
\sphinxAtStartPar
Mandatory
&
\sphinxAtStartPar
project
&
\sphinxAtStartPar
NONE
&
\sphinxAtStartPar
NONE
\\
\hline
\end{tabulary}
\par
\sphinxattableend\end{savenotes}


\subsection{Table: \sphinxstyleliteralintitle{\sphinxupquote{processing}}}
\label{\detokenize{schema_documentation:table-processing}}

\subsubsection{Description:}
\label{\detokenize{schema_documentation:id7}}

\subsubsection{Columns:}
\label{\detokenize{schema_documentation:id8}}

\begin{savenotes}\sphinxattablestart
\centering
\begin{tabulary}{\linewidth}[t]{|T|T|T|T|T|T|T|}
\hline
\sphinxstyletheadfamily 
\sphinxAtStartPar
Column
&\sphinxstyletheadfamily 
\sphinxAtStartPar
Description
&\sphinxstyletheadfamily 
\sphinxAtStartPar
Multiplicity
&\sphinxstyletheadfamily 
\sphinxAtStartPar
Obligation
&\sphinxstyletheadfamily 
\sphinxAtStartPar
Metadata level
&\sphinxstyletheadfamily 
\sphinxAtStartPar
Dublin Core
compatibility
&\sphinxstyletheadfamily 
\sphinxAtStartPar
INSPIRE
compatibility
\\
\hline
\sphinxAtStartPar
software\_name
&&&
\sphinxAtStartPar
Optional
&
\sphinxAtStartPar
project
&
\sphinxAtStartPar
NONE
&
\sphinxAtStartPar
NONE
\\
\hline
\sphinxAtStartPar
licence\_type
&&&
\sphinxAtStartPar
Optional
&
\sphinxAtStartPar
project
&
\sphinxAtStartPar
NONE
&
\sphinxAtStartPar
NONE
\\
\hline
\sphinxAtStartPar
DOI\_software
&&&
\sphinxAtStartPar
Optional
&
\sphinxAtStartPar
project
&
\sphinxAtStartPar
NONE
&
\sphinxAtStartPar
NONE
\\
\hline
\sphinxAtStartPar
notebook\_filename
&&&
\sphinxAtStartPar
Optional
&
\sphinxAtStartPar
project
&
\sphinxAtStartPar
NONE
&
\sphinxAtStartPar
NONE
\\
\hline
\sphinxAtStartPar
notebook\_purpose
&&&
\sphinxAtStartPar
Optional
&
\sphinxAtStartPar
project
&
\sphinxAtStartPar
NONE
&
\sphinxAtStartPar
NONE
\\
\hline
\sphinxAtStartPar
data\_repo\_url
&&&
\sphinxAtStartPar
Optional
&
\sphinxAtStartPar
project
&
\sphinxAtStartPar
NONE
&
\sphinxAtStartPar
NONE
\\
\hline
\sphinxAtStartPar
data\_licence
&&&
\sphinxAtStartPar
Optional
&
\sphinxAtStartPar
project
&
\sphinxAtStartPar
NONE
&
\sphinxAtStartPar
NONE
\\
\hline
\end{tabulary}
\par
\sphinxattableend\end{savenotes}


\subsection{Table: \sphinxstyleliteralintitle{\sphinxupquote{prospection}}}
\label{\detokenize{schema_documentation:table-prospection}}

\subsubsection{Description:}
\label{\detokenize{schema_documentation:id9}}

\subsubsection{Columns:}
\label{\detokenize{schema_documentation:id10}}

\begin{savenotes}\sphinxattablestart
\centering
\begin{tabulary}{\linewidth}[t]{|T|T|T|T|T|T|T|}
\hline
\sphinxstyletheadfamily 
\sphinxAtStartPar
Column
&\sphinxstyletheadfamily 
\sphinxAtStartPar
Description
&\sphinxstyletheadfamily 
\sphinxAtStartPar
Multiplicity
&\sphinxstyletheadfamily 
\sphinxAtStartPar
Obligation
&\sphinxstyletheadfamily 
\sphinxAtStartPar
Metadata level
&\sphinxstyletheadfamily 
\sphinxAtStartPar
Dublin Core
compatibility
&\sphinxstyletheadfamily 
\sphinxAtStartPar
INSPIRE
compatibility
\\
\hline
\sphinxAtStartPar
datep
&&&
\sphinxAtStartPar
Optional
&
\sphinxAtStartPar
project
&
\sphinxAtStartPar
NONE
&
\sphinxAtStartPar
NONE
\\
\hline
\sphinxAtStartPar
lat
&&&
\sphinxAtStartPar
Optional
&
\sphinxAtStartPar
project
&
\sphinxAtStartPar
NONE
&
\sphinxAtStartPar
NONE
\\
\hline
\sphinxAtStartPar
longitude
&&&
\sphinxAtStartPar
Optional
&
\sphinxAtStartPar
project
&
\sphinxAtStartPar
NONE
&
\sphinxAtStartPar
NONE
\\
\hline
\sphinxAtStartPar
method
&&&
\sphinxAtStartPar
Optional
&
\sphinxAtStartPar
project
&
\sphinxAtStartPar
NONE
&
\sphinxAtStartPar
NONE
\\
\hline
\sphinxAtStartPar
spatial\_scale
&&&
\sphinxAtStartPar
Optional
&
\sphinxAtStartPar
project
&
\sphinxAtStartPar
NONE
&
\sphinxAtStartPar
NONE
\\
\hline
\sphinxAtStartPar
bound\_cond
&&&
\sphinxAtStartPar
Optional
&
\sphinxAtStartPar
project
&
\sphinxAtStartPar
NONE
&
\sphinxAtStartPar
NONE
\\
\hline
\sphinxAtStartPar
temperature
&&&
\sphinxAtStartPar
Optional
&
\sphinxAtStartPar
project
&
\sphinxAtStartPar
NONE
&
\sphinxAtStartPar
NONE
\\
\hline
\sphinxAtStartPar
temporal\_scale
&&&
\sphinxAtStartPar
Optional
&
\sphinxAtStartPar
project
&
\sphinxAtStartPar
NONE
&
\sphinxAtStartPar
NONE
\\
\hline
\sphinxAtStartPar
instrument
&&&
\sphinxAtStartPar
Optional
&
\sphinxAtStartPar
project
&
\sphinxAtStartPar
NONE
&
\sphinxAtStartPar
NONE
\\
\hline
\sphinxAtStartPar
dimension
&&&
\sphinxAtStartPar
Optional
&
\sphinxAtStartPar
project
&
\sphinxAtStartPar
NONE
&
\sphinxAtStartPar
NONE
\\
\hline
\sphinxAtStartPar
permanent\_setup
&&&
\sphinxAtStartPar
Optional
&
\sphinxAtStartPar
project
&
\sphinxAtStartPar
NONE
&
\sphinxAtStartPar
NONE
\\
\hline
\end{tabulary}
\par
\sphinxattableend\end{savenotes}


\section{File\sphinxhyphen{}level geophysical metadata}
\label{\detokenize{schema_documentation:file-level-geophysical-metadata}}
\sphinxAtStartPar
The \sphinxstylestrong{File\sphinxhyphen{}level Metadata} level collector is a gui designed to help with the initial preparation of one geophysical dataset. Starting from one or multiple input directories, a cleanly structured output directory is generated (without deleting any input files).
Generate suitable metadata from user input and write this metadata into the directory structure, making it ready for further distribution.

\sphinxAtStartPar
\sphinxhref{https://github.com/m-weigand/geometadp.git}{geophysical Metadata Management using a Juypter Notebook}

\sphinxAtStartPar
Tables below are extracted from the database and generated by MySQL Workbench Model Documentation v1.0.0 \sphinxhyphen{} Copyright (c)
2015 Hieu Le


\subsection{Table: \sphinxstyleliteralintitle{\sphinxupquote{report}}}
\label{\detokenize{schema_documentation:table-report}}

\subsubsection{Description:}
\label{\detokenize{schema_documentation:id11}}
\sphinxAtStartPar
Metadata describing general information about the contribution.

\sphinxAtStartPar
This table has overlapping field with project level metadata (see column
CAGS Metadata level)


\subsubsection{Columns:}
\label{\detokenize{schema_documentation:id12}}

\begin{savenotes}\sphinxattablestart
\centering
\begin{tabulary}{\linewidth}[t]{|T|T|T|T|T|T|T|T|T|T|}
\hline
\sphinxstyletheadfamily 
\sphinxAtStartPar
Column
&\sphinxstyletheadfamily 
\sphinxAtStartPar
Description
&\sphinxstyletheadfamily 
\sphinxAtStartPar
Multiplicity
&\sphinxstyletheadfamily 
\sphinxAtStartPar
Obligation
&\sphinxstyletheadfamily 
\sphinxAtStartPar
CAGS Metadata level
&\sphinxstyletheadfamily 
\sphinxAtStartPar
Dublin Core (ArchSearch)
&\sphinxstyletheadfamily 
\sphinxAtStartPar
INSPIRE Directive
&\sphinxstyletheadfamily 
\sphinxAtStartPar
Data type
&\sphinxstyletheadfamily 
\sphinxAtStartPar
Attributes
&\sphinxstyletheadfamily 
\sphinxAtStartPar
Default
\\
\hline
\sphinxAtStartPar
Report\_title
&
\sphinxAtStartPar
Sort title description of the dataset
&
\sphinxAtStartPar
{[}1{]}
&
\sphinxAtStartPar
Mandatory
&
\sphinxAtStartPar
File level
&&&&
\sphinxAtStartPar
Unique
&\\
\hline
\sphinxAtStartPar
Report\_author
&
\sphinxAtStartPar
Reporting authors names
&
\sphinxAtStartPar
{[}1,n{]}
&
\sphinxAtStartPar
Mandatory
&
\sphinxAtStartPar
File \& project level
&&&&
\sphinxAtStartPar
Unique
&\\
\hline
\end{tabulary}
\par
\sphinxattableend\end{savenotes}


\subsection{Table: \sphinxstyleliteralintitle{\sphinxupquote{survey}}}
\label{\detokenize{schema_documentation:table-survey}}

\subsubsection{Description:}
\label{\detokenize{schema_documentation:description-1}}\label{\detokenize{schema_documentation:id13}}
\sphinxAtStartPar
Metadata describing one to multiple survey(s).

\sphinxAtStartPar
The survey table is inspired from \textasciigrave{}Archaeology Data Service / Digital
Antiquity

\sphinxAtStartPar
Guides to Good Practice
\textless{}\sphinxurl{https://guides.archaeologydataservice.ac.uk/g2gp/Geophysics\_6}\textgreater{}\textasciigrave{}\_

\sphinxAtStartPar
For multiple acquisitions the number n must be unchanged between the
different fields. For example, if date of time of measurement contains 2
values, the electrode configuration must contain n columns describing
the configuration used.


\subsubsection{Columns:}
\label{\detokenize{schema_documentation:columns-1}}\label{\detokenize{schema_documentation:id14}}

\begin{savenotes}\sphinxattablestart
\centering
\begin{tabulary}{\linewidth}[t]{|T|T|T|T|T|T|T|T|T|T|}
\hline
\sphinxstyletheadfamily 
\sphinxAtStartPar
Column
&\sphinxstyletheadfamily 
\sphinxAtStartPar
Description
&\sphinxstyletheadfamily 
\sphinxAtStartPar
Multiplicity
&\sphinxstyletheadfamily 
\sphinxAtStartPar
Obligation
&\sphinxstyletheadfamily 
\sphinxAtStartPar
CAGS

\sphinxAtStartPar
Metadata   level
&\sphinxstyletheadfamily 
\sphinxAtStartPar
Dublin   Core (ArchSearch)
&\sphinxstyletheadfamily 
\sphinxAtStartPar
INSPIRE   Directive
&\sphinxstyletheadfamily 
\sphinxAtStartPar
Data   type
&\sphinxstyletheadfamily 
\sphinxAtStartPar
Attributes
&\sphinxstyletheadfamily 
\sphinxAtStartPar
Default
\\
\hline
\sphinxAtStartPar
Survey\_type
&
\sphinxAtStartPar
Choose acronyms describing the survey   type (refer to CAGS glossary)
&
\sphinxAtStartPar
{[}1,n{]}
&
\sphinxAtStartPar
Mandatory
&
\sphinxAtStartPar
File level
&
\sphinxAtStartPar
Resource Type
&&&
\sphinxAtStartPar
Unique
&\\
\hline
\sphinxAtStartPar
Instruments
&
\sphinxAtStartPar
Name of the instrument(s)
&
\sphinxAtStartPar
{[}1,n{]}
&
\sphinxAtStartPar
Mandatory
&
\sphinxAtStartPar
File level
&&&&&
\sphinxAtStartPar
NULL
\\
\hline
\sphinxAtStartPar
Choice\_survey
&
\sphinxAtStartPar
Text explaining shortly the motivation of   using the method(s)
&
\sphinxAtStartPar
{[}1,n{]}
&
\sphinxAtStartPar
Optional
&
\sphinxAtStartPar
File level
&&&&&
\sphinxAtStartPar
NULL
\\
\hline
\sphinxAtStartPar
Area
&
\sphinxAtStartPar
Total surface investigated (m2)
&
\sphinxAtStartPar
{[}1,n{]}
&
\sphinxAtStartPar
Optional
&
\sphinxAtStartPar
File level
&&&&&
\sphinxAtStartPar
NULL
\\
\hline
\sphinxAtStartPar
Add\_remarks
&
\sphinxAtStartPar
Free text for additional remarks
&
\sphinxAtStartPar
{[}1{]}
&
\sphinxAtStartPar
Optional
&
\sphinxAtStartPar
File level
&&&&&
\sphinxAtStartPar
NULL
\\
\hline
\end{tabulary}
\par
\sphinxattableend\end{savenotes}


\subsection{Table: \sphinxstyleliteralintitle{\sphinxupquote{ERT metadata}}}
\label{\detokenize{schema_documentation:table-ert-metadata}}

\subsubsection{Description:}
\label{\detokenize{schema_documentation:description-2}}\label{\detokenize{schema_documentation:id15}}
\sphinxAtStartPar
Metadata describing one to multiple (n) ERT surveys.

\sphinxAtStartPar
For multiple acquisitions the number n must be unchanged between the
different fields. For example, if date of time of measurement contains 2
values, the electrode configuration must contain n columns describing
the configuration used.


\subsubsection{Columns:}
\label{\detokenize{schema_documentation:columns-3}}\label{\detokenize{schema_documentation:id16}}

\begin{savenotes}\sphinxattablestart
\centering
\begin{tabulary}{\linewidth}[t]{|T|T|T|T|T|T|T|T|T|T|}
\hline
\sphinxstyletheadfamily 
\sphinxAtStartPar
Column
&\sphinxstyletheadfamily 
\sphinxAtStartPar
Description
&\sphinxstyletheadfamily 
\sphinxAtStartPar
Multiplicity
&\sphinxstyletheadfamily 
\sphinxAtStartPar
Obligation
&\sphinxstyletheadfamily 
\sphinxAtStartPar
CAGS

\sphinxAtStartPar
Metadata level
&\sphinxstyletheadfamily 
\sphinxAtStartPar
Dublin Core (ArchSearch)
&\sphinxstyletheadfamily 
\sphinxAtStartPar
INSPIRE Directive
&\sphinxstyletheadfamily 
\sphinxAtStartPar
Data type
&\sphinxstyletheadfamily 
\sphinxAtStartPar
Attributes
&\sphinxstyletheadfamily 
\sphinxAtStartPar
Default
\\
\hline
\sphinxAtStartPar
Date\_measure
&
\sphinxAtStartPar
Date(s) of the measurement (dd/mm/aaaa)
&
\sphinxAtStartPar
{[}n{]}
&
\sphinxAtStartPar
Mandatory
&
\sphinxAtStartPar
File level
&&&&
\sphinxAtStartPar
Unique
&\\
\hline
\sphinxAtStartPar
Time\_measure
&
\sphinxAtStartPar
Time(s) of the measurement (hh:mm)
&
\sphinxAtStartPar
{[}m,n{]}
&
\sphinxAtStartPar
Mandatory
&
\sphinxAtStartPar
File level
&&&&&
\sphinxAtStartPar
NULL
\\
\hline
\sphinxAtStartPar
Elec\_conf
&
\sphinxAtStartPar
Electrode   configuration
&
\sphinxAtStartPar
{[}m,n{]}
&
\sphinxAtStartPar
Optional
&
\sphinxAtStartPar
File   level
&&&&&
\sphinxAtStartPar
NULL
\\
\hline
\sphinxAtStartPar
Elec\_spacing
&
\sphinxAtStartPar
Electrode spacing
&
\sphinxAtStartPar
{[}m,n{]}
&
\sphinxAtStartPar
Optional
&
\sphinxAtStartPar
File level
&&&&&
\sphinxAtStartPar
NULL
\\
\hline
\end{tabulary}
\par
\sphinxattableend\end{savenotes}


\subsection{Table: \sphinxstyleliteralintitle{\sphinxupquote{EM metadata}}}
\label{\detokenize{schema_documentation:table-em-metadata}}

\subsubsection{Description:}
\label{\detokenize{schema_documentation:description-3}}\label{\detokenize{schema_documentation:id17}}
\sphinxAtStartPar
Metadata describing one to multiple (n) EM surveys.

\sphinxAtStartPar
For multiple acquisitions the number n must be unchanged between the
different fields. For example, if date of measurements contains 2
values, the coil configuration must contain n columns and m lines
describing the coil configuration used.


\subsubsection{Columns:}
\label{\detokenize{schema_documentation:columns-4}}\label{\detokenize{schema_documentation:id18}}

\begin{savenotes}\sphinxattablestart
\centering
\begin{tabulary}{\linewidth}[t]{|T|T|T|T|T|T|T|T|T|T|}
\hline
\sphinxstyletheadfamily 
\sphinxAtStartPar
Column
&\sphinxstyletheadfamily 
\sphinxAtStartPar
Description
&\sphinxstyletheadfamily 
\sphinxAtStartPar
Multiplicity
&\sphinxstyletheadfamily 
\sphinxAtStartPar
Obligation
&\sphinxstyletheadfamily 
\sphinxAtStartPar
CAGS

\sphinxAtStartPar
Metadata level
&\sphinxstyletheadfamily 
\sphinxAtStartPar
Dublin Core (ArchSearch)
&\sphinxstyletheadfamily 
\sphinxAtStartPar
INSPIRE Directive
&\sphinxstyletheadfamily 
\sphinxAtStartPar
Data type
&\sphinxstyletheadfamily 
\sphinxAtStartPar
Attributes
&\sphinxstyletheadfamily 
\sphinxAtStartPar
Default
\\
\hline
\sphinxAtStartPar
Date\_measure
&
\sphinxAtStartPar
Date(s) of the measurement (dd/mm/aaaa)
&
\sphinxAtStartPar
{[}1,n{]}
&
\sphinxAtStartPar
Mandatory
&
\sphinxAtStartPar
File level
&&&&
\sphinxAtStartPar
Unique
&\\
\hline
\sphinxAtStartPar
Coil\_conf
&
\sphinxAtStartPar
Coil configuation
&
\sphinxAtStartPar
{[}m,n{]}
&
\sphinxAtStartPar
Optional
&
\sphinxAtStartPar
File level
&&&&&
\sphinxAtStartPar
NULL
\\
\hline
\sphinxAtStartPar
Read\_interval
&
\sphinxAtStartPar
If automatic time sampling, time steps   between different reading
&
\sphinxAtStartPar
{[}m,n{]}
&
\sphinxAtStartPar
Optional
&
\sphinxAtStartPar
File level
&&&&&
\sphinxAtStartPar
NULL
\\
\hline
\end{tabulary}
\par
\sphinxattableend\end{savenotes}


\subsection{Table: \sphinxstyleliteralintitle{\sphinxupquote{data quality assessment metadata}}}
\label{\detokenize{schema_documentation:table-data-quality-assessment-metadata}}

\subsubsection{Description:}
\label{\detokenize{schema_documentation:description-4}}\label{\detokenize{schema_documentation:id19}}

\subsubsection{Columns:}
\label{\detokenize{schema_documentation:columns-5}}\label{\detokenize{schema_documentation:id20}}

\begin{savenotes}\sphinxattablestart
\centering
\begin{tabulary}{\linewidth}[t]{|T|T|T|T|T|T|T|T|T|T|}
\hline
\sphinxstyletheadfamily 
\sphinxAtStartPar
Column
&\sphinxstyletheadfamily 
\sphinxAtStartPar
Description
&\sphinxstyletheadfamily 
\sphinxAtStartPar
Multiplicity
&\sphinxstyletheadfamily 
\sphinxAtStartPar
Obligation
&\sphinxstyletheadfamily 
\sphinxAtStartPar
CAGS

\sphinxAtStartPar
Metadata level
&\sphinxstyletheadfamily 
\sphinxAtStartPar
Dublin Core (ArchSearch)
&\sphinxstyletheadfamily 
\sphinxAtStartPar
INSPIRE Directive
&\sphinxstyletheadfamily 
\sphinxAtStartPar
Data type
&\sphinxstyletheadfamily 
\sphinxAtStartPar
Attributes
&\sphinxstyletheadfamily 
\sphinxAtStartPar
Default
\\
\hline
\sphinxAtStartPar
Peer\_reviewed
&
\sphinxAtStartPar
True if the dataset has been   peer\sphinxhyphen{}reviewed
&
\sphinxAtStartPar
{[}1{]}
&
\sphinxAtStartPar
Mandatory
&
\sphinxAtStartPar
File level
&&&&
\sphinxAtStartPar
Unique
&\\
\hline
\sphinxAtStartPar
Peer\_reviewer\_contact
&
\sphinxAtStartPar
Contact of reviewer
&
\sphinxAtStartPar
{[}1,n{]}
&
\sphinxAtStartPar
Mandatory only if peer\_reviewed is True
&&&&&
\sphinxAtStartPar
Unique
&
\sphinxAtStartPar
NULL
\\
\hline
\sphinxAtStartPar
Replicate\_datasets
&
\sphinxAtStartPar
Number of replicates datasets
&
\sphinxAtStartPar
{[}1,n{]}
&
\sphinxAtStartPar
Optional
&&&&&&
\sphinxAtStartPar
NULL
\\
\hline
\sphinxAtStartPar
Comparison\_ref\_data
&
\sphinxAtStartPar
The dataset has been compared with   reference datasets
&
\sphinxAtStartPar
{[}1,n{]}
&
\sphinxAtStartPar
Optional
&&&&&&
\sphinxAtStartPar
NULL
\\
\hline
\sphinxAtStartPar
Ref\_data
&
\sphinxAtStartPar
DOI of reference dataset
&
\sphinxAtStartPar
{[}1,n{]}
&
\sphinxAtStartPar
Optional
&&&&&&
\sphinxAtStartPar
NULL
\\
\hline
\end{tabulary}
\par
\sphinxattableend\end{savenotes}


\subsection{Table: \sphinxstyleliteralintitle{\sphinxupquote{sampling}}}
\label{\detokenize{schema_documentation:table-sampling}}\label{\detokenize{schema_documentation:section-1}}

\subsubsection{Description:}
\label{\detokenize{schema_documentation:description-5}}\label{\detokenize{schema_documentation:id21}}

\subsubsection{Columns:}
\label{\detokenize{schema_documentation:columns-6}}\label{\detokenize{schema_documentation:id22}}\phantomsection\label{\detokenize{schema_documentation:section-2}}

\chapter{Contributors}
\label{\detokenize{contibutors:contributors}}\label{\detokenize{contibutors::doc}}\begin{itemize}
\item {} 
\sphinxAtStartPar
Guillaume Blanchy, University of Lancaster

\end{itemize}


\chapter{Terms of use}
\label{\detokenize{terms_of_use:terms-of-use}}\label{\detokenize{terms_of_use::doc}}
\sphinxAtStartPar
The agrigeophysical catalog is made available under the \sphinxhref{https://opendatacommons.org/licenses/by/1.0/}{Open Data Commons Attribution License (ODC\sphinxhyphen{}By) v1.0}


\chapter{Indices and tables}
\label{\detokenize{index:indices-and-tables}}\begin{itemize}
\item {} 
\sphinxAtStartPar
\DUrole{xref,std,std-ref}{genindex}

\item {} 
\sphinxAtStartPar
\DUrole{xref,std,std-ref}{modindex}

\item {} 
\sphinxAtStartPar
\DUrole{xref,std,std-ref}{search}

\end{itemize}

\begin{sphinxthebibliography}{PER2020}
\bibitem[PER2020]{glossary:per2020}
\sphinxAtStartPar
Peruzzo, L., Chou, C., Wu, Y. et al. Imaging of plant current pathways for non\sphinxhyphen{}invasive root Phenotyping using a newly developed electrical current source density approach. Plant Soil 450, 567\textendash{}584 (2020). \sphinxurl{https://doi.org/10.1007/s11104-020-04529-w}
\bibitem[ENR+20]{glossary:id6}
\sphinxAtStartPar
Solomon Ehosioke, Frédéric Nguyen, Sathyanarayan Rao, Thomas Kremer, Edmundo Placencia‐Gomez, Johan Alexander Huisman, Andreas Kemna, Mathieu Javaux, and Sarah Garré. Sensing the electrical properties of roots: A review. \sphinxstyleemphasis{Vadose Zone Journal}, 19(1):e20082, 2020. \_eprint: https://acsess.onlinelibrary.wiley.com/doi/pdf/10.1002/vzj2.20082. URL: \sphinxurl{https://acsess.onlinelibrary.wiley.com/doi/abs/10.1002/vzj2.20082} (visited on 2021\sphinxhyphen{}01\sphinxhyphen{}04), \sphinxhref{https://doi.org/https://doi.org/10.1002/vzj2.20082}{doi:https://doi.org/10.1002/vzj2.20082}.
\end{sphinxthebibliography}



\renewcommand{\indexname}{Index}
\printindex
\end{document}